\section{Conclusion}
\label{sec:conclusion}

This paper set out to test, under a falsification-first protocol, whether
a maximum-likelihood-constrained multivariate Hawkes process generalizes
better than an unconstrained gradient-boosted-tree model at forecasting
L1 order-book queue depletion, in a setting deliberately constructed so
that the parametric model's intensity is a strict linear subspace of the
tree model's own feature space -- making the comparison one of
generalization under constraint, not of which model can see more
information. Across every axis tested -- the full sample-size learning
curve from 2 hours to 31 days of training data, a formal 31-day rolling
Giacomini--White test, a calibration-budget robustness sweep up to
50$\times$ the original Hawkes calibration sample, an 18-specification
Romano--Wolf FWER-controlled grid, and independently tuned LightGBM
hyperparameters for both models -- the free tree model $M_1$ beats the
parametric Hawkes model $M_2$, by a wide and statistically decisive
margin, with no crossover sample size at which the parametric model
overtakes it. The single standing hypothesis that could have explained
this away as a methodological artifact in $M_1$'s favor -- that $M_1$'s
hyperparameters were undertuned relative to $M_2$'s -- was tested
directly and found backwards: proper tuning narrows the gap by
benefiting $M_2$ more than $M_1$, not the reverse.

At the same time, the project did not simply confirm a null and stop.
$M_3$, the one model in the ladder with access to genuinely non-linear
information ($M_1$'s raw feature bank does not already contain) -- an
hourly Hawkes-derived spectral-radius regime indicator and linearized
trade-size marks -- showed no ranking-quality (PR-AUC) advantage over
$M_1$ on the 500\,ms Primary Endpoint, exactly as the hourly regime
feature's own coarser timescale would predict. Yet a formal test on a
different, calibration-sensitive metric (log-loss, under the same
31-day rolling protocol that established the $M_1$-vs-$M_2$ result)
revealed a small but highly significant improvement, which an ablation
traced cleanly to the regime feature alone, with the trade-size marks
contributing nothing and mildly diluting the combined effect. This
result -- real on one metric, absent on another, and not large enough to
clear the pre-registered PR-AUC magnitude threshold under either reading
-- is reported as exactly that: a modest, metric-specific finding, not
rounded up into a headline reversal of the paper's primary conclusion,
consistent with the project's stated commitment to reporting what the
data says rather than what would make the best story.

The cross-asset transfer ladder adds a further, independent line of
evidence on the value of the parametric compression itself, separate
from the sample-efficiency question the learning curve addresses. A
branching-ratio matrix fit purely on BTC, transferred with only a local
baseline intensity and marks re-estimated per target, recovers 81.7\%--89.3\%
of each of three structurally different assets' (ETH, SOL, DOGE) own
free-model advantage over the $M_0$ baseline -- with no evidence of
degradation as tick size and event mix diverge further from BTC's own.
This is a genuine, if narrower, generalization advantage for the
parametric representation: not better absolute forecasting accuracy than
a locally-trained free model, but a substantial fraction of that
accuracy transported across instruments at essentially zero re-training
cost, which a black-box tree model has no comparable native mechanism to
achieve.

Taken as a whole, the evidence assembled here falsifies the hypothesis
that MLE-constrained parametric compression outperforms unconstrained
tree learning for L1 queue-depletion \emph{ranking} quality at any tested
sample size, training-data diversity, or calibration budget on this
instrument and window -- a clean, negative, and well-powered result in
its own right. It simultaneously identifies two narrower channels in
which the parametric structure retains real value: a calibration-only
predictive edge from a single Hawkes-derived regime summary, and a
substantial, non-degrading transfer of predictive structure across
instruments that a purely data-driven free model cannot replicate
without retraining. Future work most directly motivated by this paper's
own findings and limitations
(Section~\ref{sec:limitations}) includes replicating the primary
comparison over additional calendar windows and exchanges to establish
temporal and cross-venue external validity, running the specified
difference-in-differences placebo control against $M_3$'s
now-established calibration effect rather than its originally tested
null PR-AUC result, and extending the transfer ladder with a formal
significance test and additional BTC source windows now that its
disclosed simplifications have been resolved.
